\documentclass[../DoAn.tex]{subfiles}
\begin{document}

Chương 1 giới thiệu bối cảnh, động lực và phạm vi của đồ án. Từ bài toán điều hướng
đa tác nhân trong game, chương này xác định mục tiêu cụ thể, đề xuất định hướng giải
pháp và trình bày tổng quan cấu trúc báo cáo.

% ============================================================
\section{Đặt vấn đề}
\label{section:1.1}

Trong các trò chơi điện tử chiến thuật thời gian thực (Real-Time Strategy -- RTS),
điều hướng đồng thời nhiều nhân vật AI đến mục tiêu là một thách thức kỹ thuật cốt lõi.
Khi số lượng nhân vật AI tăng lên, xác suất xảy ra va chạm, tình trạng deadlock và
đường đi không hiệu quả cũng tăng theo. Đây chính là bài toán MAPF -- Multi-Agent
Pathfinding -- được nghiên cứu rộng rãi trong lĩnh vực trí tuệ nhân tạo và robot học
\cite{stern2019mapf}.

MAPF yêu cầu lập kế hoạch đồng thời cho $k$ agent trên đồ thị lưới để đảm bảo mỗi
agent đến đích trong thời gian tối thiểu mà không va chạm nhau tại bất kỳ bước thời
gian nào. Trong môi trường game thời gian thực, bài toán này đặt ra thêm ràng buộc
về thời gian tính toán: thuật toán phải phản hồi trong vài mili-giây và phải tái lập
kế hoạch liên tục khi môi trường thay đổi.

Các thuật toán tìm đường agent đơn như A* \cite{hart1968astar} được sử dụng phổ biến
trong game nhờ độ đơn giản và hiệu quả, nhưng khi áp dụng độc lập cho nhiều agent,
chúng không đảm bảo tránh va chạm. Các thuật toán MAPF tối ưu như CBS
\cite{sharon2015cbs} đảm bảo tính tối ưu toàn cục nhưng có độ phức tạp tính toán
tăng theo hàm mũ với số agent, không phù hợp cho game thời gian thực. PIBT (Priority
Inheritance with Backtracking) \cite{okumura2022pibt} là một hướng tiếp cận cân bằng:
chạy theo bước thời gian, hoàn chỉnh và hiệu quả cho môi trường thời gian thực, song
chấp nhận thay thế lời giải tối ưu toàn cục bằng lời giải khả thi đủ tốt theo từng
bước thời gian để đạt tốc độ phản hồi thời gian thực.

PIBT trong đồ án được tham khảo từ cuộc thi điều phối robot The League of Robot Runners
(LoRR) \cite{lorr2024} do Amazon Robotics tài trợ, nơi nhóm Team No Man's Sky không chỉ
giành giải nhất ở một hạng mục mà còn đứng nhất cả ba bảng đấu -- Combined, Planner và
Scheduler -- đồng thời đoạt giải Line Honours (nhất toàn đoàn với số lời giải tốt nhất
nhiều nhất) trong Main Round 2024. Tuy nhiên, các cài đặt đoạt giải đều ở dạng máy chủ
C++ chạy benchmark, không thể chơi hay tương tác trực tiếp. Đồ án kế thừa ý tưởng và
thuật toán của Team No Man's Sky rồi đưa vào Unity, để người chơi có thể trực tiếp trải
nghiệm và quan sát hành vi phối hợp đa agent một cách trực quan ngay trong game.

Thực tế trong ngành game, phần lớn các tựa game RTS và tower-defense hiện tại sử dụng
thuật toán tìm đường đơn giản hoặc Unity NavMesh mà không giải quyết bài toán va chạm
đa agent ở lớp lập kế hoạch. Điều này dẫn đến các hiện tượng nhân vật xếp hàng chờ
hoặc chen lấn không tự nhiên, làm giảm chất lượng trải nghiệm người chơi.

Đồ án này xuất phát từ nhu cầu nghiên cứu và so sánh thực nghiệm ba cấp độ thuật toán
MAPF -- từ baseline A* đến PIBT phối hợp đa agent và PIBT tích hợp máy chủ C++ bên
ngoài -- trong bối cảnh game bắn xe tăng 2D nhiều người chơi xây dựng trên Unity Engine.
Kết quả sẽ cung cấp bằng chứng thực nghiệm về ưu nhược điểm của từng hướng tiếp cận
trong môi trường game thực tế.

% ============================================================
\section{Mục tiêu và phạm vi đề tài}
\label{section:1.2}

Các sản phẩm và nghiên cứu liên quan hiện tại có thể phân thành ba nhóm. Nhóm thứ nhất
là các game thương mại sử dụng tìm đường đơn giản: phần lớn game RTS và tower-defense
dùng A* hoặc Unity NavMesh cho từng agent độc lập, không giải quyết va chạm đa agent ở
lớp lập kế hoạch, dẫn đến hiện tượng nhân vật xếp hàng hoặc chen lấn không tự nhiên.
Nhóm thứ hai là các nghiên cứu MAPF học thuật như CBS, EECBS hay PIBT, được nghiên cứu
sâu trong cộng đồng trí tuệ nhân tạo nhưng thường chỉ được đánh giá trên benchmark thuần
túy mà không tích hợp trực tiếp vào engine game thương mại. Nhóm thứ ba là một số nghiên
cứu thử tích hợp MAPF vào game \cite{de2013cooperative}, song thường chỉ chạy trên môi
trường giả lập đơn giản, thiếu một engine vật lý đầy đủ.

Từ ba nhóm trên, có thể thấy khoảng trống hiện tại là thiếu một so sánh thực nghiệm đầy
đủ giữa các cấp độ MAPF trong môi trường game có vật lý thực tế (Unity physics) trên tập
bản đồ benchmark chuẩn. Đây chính là khoảng trống mà đồ án hướng tới lấp đầy.

Trên cơ sở đó, đồ án đặt ra mục tiêu xây dựng một game bắn xe tăng 2D nhiều người chơi
trên Unity sử dụng bản đồ chuẩn MovingAI MAPF benchmark, đồng thời triển khai và tích hợp
ba cấp độ thuật toán MAPF gồm A* làm baseline, PIBT~C\# và PIBT~C++/TCP. Để so sánh ba
thuật toán một cách khách quan, đồ án xây dựng thêm một hệ thống backtest tự động đo lường
định lượng trên năm bản đồ, từ đó phân tích kết quả và rút ra nhận xét về khả năng ứng
dụng của từng thuật toán trong game thời gian thực.

Về phạm vi, đồ án giới hạn ở game 2D góc nhìn từ trên xuống, môi trường lưới rời rạc
(discrete grid), tối đa sáu agent AI trong kịch bản thực nghiệm, và không xét tối ưu hóa
đa mục tiêu phức tạp.

% ============================================================
\section{Định hướng giải pháp}
\label{section:1.3}

Từ phân tích ở Mục~\ref{section:1.2}, đồ án đề xuất ba bước triển khai:

Thứ nhất, xây dựng hạ tầng bản đồ thống nhất dựa trên đọc tệp \texttt{.map} MovingAI
và dựng lưới ô tại thời điểm chạy trong Unity, thay vì bake cứng vào scene. Cách này
cho phép thử nghiệm trên nhiều bản đồ mà không cần thay đổi code.

Thứ hai, triển khai tuần tự ba thuật toán MAPF chia sẻ cùng tầng điều khiển vật lý
(\texttt{TankController}), đảm bảo hành vi vật lý nhất quán khi so sánh thuật toán.
A* đóng vai trò baseline; PIBT~C\# thêm phối hợp đa agent ngay trong Unity; PIBT~C++/TCP
mở rộng tích hợp máy chủ ngoài.

Thứ ba, xây dựng hệ thống backtest tự động chạy kịch bản, thu thập chỉ số và xuất CSV,
cho phép so sánh định lượng trên tập bản đồ lớn mà không cần can thiệp thủ công. Thêm
vào đó, module chướng ngại động (spawn/despawn thùng cản đường trực tiếp trên path
của agent) dùng để stress-test khả năng tái hoạch định của từng thuật toán.

Ngoài ba bước trên, thiết kế hệ thống còn hướng tới khả năng mở rộng quy mô về sau trên
ba trục: tăng số lượng agent AI, tăng số người chơi đồng thời, và mở rộng hạ tầng máy
chủ (đưa PIBT server lên cloud, cân bằng tải). Các hướng mở rộng này được trình bày chi
tiết ở Chương~\ref{chapter:conclusion}.

Đóng góp chính: hệ thống tích hợp MAPF đầy đủ trong Unity với ba cấp độ thuật toán,
cơ sở hạ tầng backtest định lượng, và bộ kết quả thực nghiệm trên benchmark chuẩn.

% ============================================================
\section{Bố cục đồ án}
\label{section:1.4}

Phần còn lại của báo cáo đồ án tốt nghiệp này được tổ chức như sau.

Chương~2 trình bày khảo sát hiện trạng các giải pháp điều hướng đa agent trong game,
phân tích yêu cầu chức năng và phi chức năng của hệ thống, đặc tả các use case chính
thông qua biểu đồ use case và đặc tả luồng sự kiện.

Chương~3 trình bày nền tảng lý thuyết và công nghệ được sử dụng: bài toán MAPF và
cách biểu diễn, thuật toán A* và heuristic Manhattan, thuật toán PIBT và cơ chế kế
thừa ưu tiên có hoàn tác, cùng với lý do lựa chọn Unity Engine và giao thức TCP cho
tích hợp máy chủ ngoài.

Chương~4 trình bày phân tích thiết kế kiến trúc hệ thống theo mô hình phân tầng,
thiết kế chi tiết các lớp MAPF AI, quá trình xây dựng ứng dụng với thống kê mã nguồn,
và đánh giá thực nghiệm với số liệu backtest trên 90 kịch bản tĩnh và 90 kịch bản
có chướng ngại động.

Chương~5 trình bày năm đóng góp kỹ thuật nổi bật: hệ thống đọc bản đồ động, triển khai
A* thời gian thực, tích hợp PIBT phối hợp đa agent, cầu nối TCP đến máy chủ C++, và
hệ thống backtest tự động với chướng ngại động.

Chương~6 tổng kết kết quả đã đạt được, so sánh với mục tiêu đề ra, và định hướng
phát triển trong tương lai.

\textbf{Kết chương~1:} Chương này đã xác định bài toán MAPF trong game thời gian thực,
nêu rõ khoảng trống nghiên cứu và đề xuất hướng giải quyết bằng cách tích hợp và so
sánh ba thuật toán (A*, PIBT~C\#, PIBT~C++/TCP) trong môi trường Unity. Các chương tiếp
theo sẽ trình bày chi tiết quá trình phân tích, thiết kế, triển khai và đánh giá hệ thống.

\end{document}
